                \begin{tikzpicture}[xscale=1.0, yscale=1.0]
                    \useasboundingbox (-0.35, -1.25) rectangle (11.55, 4.35);

                    \tikzset{cadre/.style={draw=gray!45, line width=0.9pt}}
                    \tikzset{coupe1/.style={gris_fonce_inria, line width=1.6pt}}
                    \tikzset{coupe2/.style={gris_fonce_inria, line width=0.8pt, opacity=0.75}}
                    \tikzset{coupe3/.style={gris_fonce_inria, line width=0.45pt, opacity=0.5}}
                    \tikzset{fissure/.style={inria-rouge, line width=2.4pt, line cap=round}}
                    \tikzset{bloc/.style={fill=inria-2024-bleu-canard!16, draw=inria-2024-bleu-canard,
                                          line width=0.5pt}}

                    % ================= gauche : hiérarchie équilibrée =================
                    \begin{scope}[xshift=0.9cm, yshift=0.55cm]
                        \fill[gray!4] (0,0) rectangle (3.6,3.6);
                        % coupes emboîtées, équilibrées
                        \draw[coupe1] (1.8,0) -- (1.8,3.6);
                        \draw[coupe2] (0,1.8) -- (1.8,1.8);
                        \draw[coupe2] (1.8,1.8) -- (3.6,1.8);
                        \foreach \x in {0.9, 2.7} { \draw[coupe3] (\x,0) -- (\x,3.6); }
                        \foreach \y in {0.9, 2.7} { \draw[coupe3] (0,\y) -- (3.6,\y); }
                        \draw[cadre] (0,0) rectangle (3.6,3.6);
                        % la fissure traverse
                        \draw[fissure] (0.25,3.25) .. controls (1.20,2.45) and (1.45,1.55) .. (2.35,0.95)
                                       .. controls (2.85,0.62) and (3.10,0.55) .. (3.35,0.30);
                        % point de rupture
                        \fill[white, draw=inria-rouge, line width=1.1pt] (1.80,1.62) circle (4.6pt);
                        \node[font=\tiny\bfseries, text=inria-rouge] at (1.80,1.62) {$\times$};
                        \node[font=\scriptsize, text=gris_fonce_inria, anchor=north, align=center] at (1.8,-0.24)
                            {\textbf{hiérarchie équilibrée}\\ce qu'exige l'efficacité};
                    \end{scope}

                    % ================= droite : coupe sémantique =================
                    \begin{scope}[xshift=6.6cm, yshift=0.55cm]
                        \fill[gray!4] (0,0) rectangle (3.6,3.6);
                        % bloc « fissure » : bande le long de la fissure
                        \fill[inria-rouge!10, draw=inria-rouge!45, line width=0.7pt]
                            plot[smooth] coordinates {(0.05,3.55) (1.10,2.60) (1.55,1.70) (2.45,1.15) (3.55,0.55)}
                            -- plot[smooth] coordinates {(3.55,0.05) (2.35,0.75) (1.15,1.35) (0.75,2.35) (0.05,3.05)}
                            -- cycle;
                        % fond découpé en cellules compactes
                        \draw[coupe3] (2.05,3.6) -- (3.6,2.05);
                        \draw[coupe3] (2.75,3.6) -- (3.6,2.75);
                        \draw[coupe3] (0,1.75) -- (1.45,0.30);
                        \draw[coupe3] (0,1.05) -- (0.80,0.25);
                        \draw[cadre] (0,0) rectangle (3.6,3.6);
                        \draw[fissure] (0.25,3.25) .. controls (1.20,2.45) and (1.45,1.55) .. (2.35,0.95)
                                       .. controls (2.85,0.62) and (3.10,0.55) .. (3.35,0.30);
                        \node[font=\scriptsize, text=gris_fonce_inria, anchor=north, align=center] at (1.8,-0.24)
                            {\textbf{coupe \{fissure, fond\}}\\la fissure tient dans un sous-arbre};
                    \end{scope}

                    % ================= verdict =================
                    \node[font=\scriptsize, align=center, text=gris_fonce_inria] at (5.65,3.30)
                        {une coupe équilibrée\\de niveau $1$\\\textbf{traverse} la fissure};
                    \node[font=\Large\bfseries, text=inria-rouge] at (5.65,1.95) {$\neq$};
                    \node[font=\scriptsize, align=center, text=gris_fonce_inria] at (5.65,0.95)
                        {efficacité\\\emph{contre}\\continuité};
                \end{tikzpicture}
